\section{Classical Graph Neural Networks as Propagation}
\label{sec:classical}

\subsection{Message passing is local propagation}
A graph convolutional network (GCN) layer~\cite{kipf2017gcn,defferrard2016chebnet}
updates node embeddings by averaging over neighbors and applying a learned linear map
and nonlinearity,
\begin{equation}
\Feat^{(\ell+1)}=\sigma\!\big(\hat{\Adj}\,\Feat^{(\ell)}\Wt^{(\ell)}\big),\qquad
\hat{\Adj}=\tilde{\Deg}^{-1/2}\tilde{\Adj}\,\tilde{\Deg}^{-1/2},
\end{equation}
where $\tilde{\Adj}=\Adj+\Id$ adds self-loops and $\tilde{\Deg}$ is its degree matrix.
The operator $\hat{\Adj}$ is a degree-one polynomial in the adjacency, so one layer
mixes one hop and $K$ stacked layers mix exactly $K$ hops. This locality is the key
structural fact about GCNs: a $K$-layer GCN simply cannot move information farther
than $K$ hops, regardless of width.

\subsection{The continuous view: diffusion}
Stacking many small averaging steps approaches a continuous diffusion. The heat
equation on a graph, $\dot{\mathbf{x}}=-\Lap\mathbf{x}$, has the closed-form solution
\begin{equation}
\mathbf{x}(t)=e^{-t\Lap}\,\mathbf{x}(0),\qquad
e^{-t\Lap}=\Eig\,\mathrm{diag}\!\big(e^{-t\lambda_k}\big)\,\Eig^{\top}.
\label{eq:heat}
\end{equation}
Because $\Lap\mathbf{1}=\mathbf{0}$, the kernel $e^{-t\Lap}$ is row-stochastic and
nonnegative: it is a genuine probability kernel that smooths a signal across the
whole graph, with a single scalar $t$ controlling how far. A diffusion-based GNN
layer~\cite{gasteiger2019diffusion} replaces $\hat{\Adj}$ by $e^{-t\Lap}$ and, unlike a
fixed $K$-hop GCN, can reach any node at the cost of a learnable time $t$. Other local
message-passing schemes~\cite{hamilton2017graphsage,velickovic2018gat} share the same
$K$-hop horizon.

\subsection{Why range matters}
The defining property of diffusion is \emph{how fast it spreads}. For an impulse
released at a node, the spatial second moment of the heat kernel grows linearly in
time,
\begin{equation}
\Var\big[e^{-t\Lap}\big]\;\sim\;t .
\label{eq:diffusive}
\end{equation}
This linear, diffusive spreading is the classical baseline against which we will
contrast the quantum walk. It already tells us what to expect from the case study:
on tasks that require long-range information, a global diffusion operator should beat
a local $K$-hop GCN, simply because it can reach the far nodes at all. The interesting
question is what, if anything, a quantum operator adds beyond reaching them.

\paragraph*{Worked example} Take a path of $N=81$ nodes and release an impulse at the
center. The heat kernel $e^{-t\Lap}$ gives a Gaussian-like profile whose spatial
variance is exactly $2t$: at $t=1$ the variance is $2$, at $t=4$ it is $8$, and at
$t=8$ it is $16$ (the steel curve in Fig.~\ref{fig:expA}). The mass reaches farther as
$t$ grows, but only as $\sqrt{t}$ in distance. A $K$-layer GCN, by contrast, places
exactly zero mass beyond hop $K$ no matter how it is trained; on this path a $3$-layer
GCN can never see node $40$ from the center. This is the concrete sense in which range,
not capacity, is the limiting resource for local operators.
